\section{Khai báo sử dụng AI}
\label{sec:ai-disclosure}

\subsection{Phạm vi sử dụng}

Nhóm sử dụng OpenAI Codex để hỗ trợ lập kế hoạch, đối chiếu báo cáo với mã
nguồn, chỉnh sửa heuristic A*, kiểm thử và trình bày kết quả. Codex đưa ra đề
xuất hoặc thực hiện thao tác theo yêu cầu; nhóm kiểm tra đề xuất về thuật toán
với luật chơi và hàm chi phí, kiểm tra phần mô tả với mã nguồn. Số liệu thực
nghiệm lấy từ chương trình đo, không lấy từ câu trả lời do AI tự tạo. Trong
quá trình rà soát, nhóm còn tham khảo các hướng dẫn \textit{peer-review},
\textit{scientific-writing} và \textit{scientific-visualization} thuộc bộ
\textit{Scientific Agent Skills}~\cite{scientificagentskills2026}.

\subsection{Các tương tác tiêu biểu}

Các đoạn dưới đây là trích đoạn nguyên văn từ yêu cầu nhóm đã gửi Codex trong
các phiên làm việc của đồ án; cách viết và lỗi chính tả trong prompt được giữ
nguyên. Chúng cho thấy mục đích sử dụng công cụ, không đại diện cho toàn bộ
lịch sử trao đổi.

\subsubsection{Lập kế hoạch báo cáo và video}

\noindent\textbf{Prompt (trích đoạn):}
\begin{quote}
\small\itshape
Tôi muốn bạn dựa vào nội dung yêu cầu của đề bài thì hãy tạo ra cho tôi một
file kế hoạch viết report + làm video
\end{quote}

Nhóm dùng tương tác này để chuyển các hạng mục của đề bài thành kế hoạch thực
hiện: phân công, mô tả thuật toán, thực nghiệm, tài liệu tham khảo và video
demo. Kế hoạch là tài liệu tổ chức công việc; nhóm vẫn đối chiếu từng hạng mục
với đề bài chính thức trước khi đưa vào bản nộp.

\subsubsection{Đối chiếu báo cáo với chương trình}

\noindent\textbf{Prompt (trích đoạn):}
\begin{quote}
\small\itshape
Đối chiếu toàn bộ Chương 1 và Chương 2 của báo cáo với implementation thực
tế.
\end{quote}

Codex được yêu cầu tìm những chỗ mô tả trạng thái, luật chơi, kiến trúc hoặc
chức năng giao diện chưa khớp với mã nguồn. Nhóm dùng kết quả rà soát để sửa
cách diễn đạt và kiểm tra lại các tệp hiện thực liên quan. Vì vậy, phần mô tả
trong báo cáo dựa trên hành vi của chương trình, không chỉ dựa trên bản thảo
do AI đề xuất.

\subsubsection{Thiết kế lại heuristic A*}

\noindent\textbf{Prompt (trích đoạn):}
\begin{quote}
\small\itshape
Vậy thì bạn chọn cho tôi một cái heuristic mới đi á đảm bảo tính admissible
á. Sửa luôn code cho tôi á.
\end{quote}

Tương tác này liên quan trực tiếp tới phần hiện thực: Codex hỗ trợ xây dựng và
chỉnh sửa heuristic để phù hợp với hàm chi phí 0/1/2 và cơ chế tách khối.
Nhóm đối chiếu kết quả với luật chơi, bộ kiểm thử và lập luận cận dưới ở
Mục~\ref{sec:algorithms}. Phép kiểm tra hữu hạn trên bộ màn chơi bổ sung cho
lập luận, không thay thế việc chứng minh điều kiện áp dụng.

\subsubsection{Chạy lại thực nghiệm}

\noindent\textbf{Prompt (trích đoạn):}
\begin{quote}
\small\itshape
Bây giờ tôi muốn bạn kiểm tra xem chúng ta có phải là source code mới nhất
chưa TH\_AI\_LABS\_1 á rồi pull cái main mới nhất về. Rồi sau đó chạy lại
benmark lại một lần nữa cho tôi á
\end{quote}

Codex hỗ trợ cập nhật mã nguồn, chạy lại benchmark và đồng bộ CSV, bảng cùng
biểu đồ trong báo cáo. Bộ số liệu cuối gồm 180 lượt chạy trên 15 màn chơi;
tập lệnh kiểm chứng đối chiếu CSV và xác nhận 60/60 lời giải hợp lệ. Nhóm dùng
những tệp kết quả này để viết phần phân tích, đồng thời nêu giới hạn về số màn
chơi, số lần lặp và phép đo bộ nhớ.

\subsection{Kiểm chứng và trách nhiệm}

Nhóm chịu trách nhiệm về lựa chọn thuật toán, tính đúng đắn của mã nguồn, cách
diễn giải số liệu, nội dung báo cáo và mức độ hoàn thành của từng thành viên.
Mã nguồn dùng cho thực nghiệm đạt 203/203 trường hợp kiểm thử. Những nội dung
Codex đề xuất chỉ được sử dụng sau khi đối chiếu với đề bài, mã nguồn hoặc dữ
liệu thực nghiệm tương ứng. Các trích đoạn trên không phải danh sách đầy đủ
các prompt.
